\documentclass[11pt]{article}
\usepackage[margin=2.2cm]{geometry}
\usepackage{fontspec}
\setmainfont{STIX Two Text}
\usepackage{unicode-math}
\setmathfont{STIX Two Math}
\usepackage{graphicx,booktabs,enumitem,hyperref,float,amsmath}\hypersetup{hidelinks}
\setlist{nosep,leftmargin=1.4em}
\setlength{\parskip}{5pt}\setlength{\parindent}{0pt}\setlength{\emergencystretch}{2em}
\title{Two-color excitation with optical pulses and a delay:\\ derivations, step by step}
\date{12 November 2026, revised 19 November 2026}
\author{}
\begin{document}
\maketitle
\vspace{-2.5em}

This document continues \texttt{two\_color\_plan.pdf}, which builds the focal field of one lens from Maxwell's
equations and the two-color signal for continuous beams. Here each step from a Gaussian pulse to the delay
results is written out, with no step skipped, and each result is compared with the stored computations. The
equations of the papers are cited where they enter: Cambaliza \& Saloma (2000) [CS00], Lim \& Saloma (2002) [LS02],
Lim \& Saloma (2003) [LS03] and Romallosa, Bantang \& Saloma (2003) [Rom03].

\section*{Symbols}
\begin{center}\small\begin{tabular}{ll}\toprule
$\lambda_1, \lambda_2$ & wavelengths of beams 1 and 2, 750 and 656 nm\\
$\omega_1, \omega_2$ & their angular frequencies, $2\pi c/\lambda$: 2.5115 and 2.8714 rad/fs\\
$\omega_e$ & angular frequency of the level, $2\pi c/(350\ \mathrm{nm}) = 5.3819$ rad/fs\\
$\tau$ & pulse width: the FWHM of the intensity $|E(t)|^2$\\
$\Delta t$ & delay of pulse 2 after pulse 1\\
$\theta$ & angle between the two lens axes: 0 collinear, $\pi/2$ crossed, $\pi$ facing\\
$n$ & refractive index of the medium; $k = n\omega/c$\\
$z_b$ & distance of a point from the focus of a beam, measured along that beam\\
$\langle\cos\vartheta\rangle$ & mean of $\cos\vartheta$ over the cone of rays, each ray weighted by its amplitude\\
$v$ & speed of a pulse along the axis near the focus, $c/(n\langle\cos\vartheta\rangle)$\\
$\mathbf{E}_\omega(P)$ & focal field of a continuous wave of frequency $\omega$ at the point $P$\\
$c$ & speed of light, 0.29979 $\mu$m/fs\\\bottomrule\end{tabular}\end{center}

\section{The Gaussian pulse}
\textbf{Step 1.} I write the field of the pulse as a carrier times a Gaussian envelope,
\begin{equation}E(t) = e^{-at^2}\,e^{-i\omega_0t},\qquad a > 0 .\end{equation}
This is Eq.~8 of [Rom03] with no chirp ($b = 0$).

\textbf{Step 2.} The intensity is the squared modulus. The carrier has modulus 1, so
\begin{equation}|E(t)|^2 = \left(e^{-at^2}\right)^2 = e^{-2at^2} .\end{equation}

\textbf{Step 3.} The FWHM $\tau$ is the full width at which the intensity is one half of its peak. The peak is at
$t = 0$ and equals 1. Setting $t = \tau/2$:
\begin{align}e^{-2a(\tau/2)^2} &= \tfrac12\\ -2a\,\frac{\tau^2}{4} &= -\ln 2\\ a &= \frac{2\ln 2}{\tau^2} .\end{align}

\textbf{Step 4.} Putting $a$ back,
\begin{equation}E(t) = \exp\!\left(-2\ln 2\,\frac{t^2}{\tau^2}\right)e^{-i\omega_0t},\qquad
|E(t)|^2 = \exp\!\left(-4\ln 2\,\frac{t^2}{\tau^2}\right) ,\label{eq:pulse}\end{equation}
which is Eq.~9 of [Rom03], where $\tau = (2\ln 2/a)^{1/2}$.

\section{The spectrum of the pulse}
\textbf{Step 1.} The spectrum is the Fourier transform of the field,
\begin{equation}\tilde E(\omega) = \int_{-\infty}^{\infty}E(t)\,e^{i\omega t}\,dt
= \int_{-\infty}^{\infty}e^{-at^2}\,e^{i(\omega-\omega_0)t}\,dt .\end{equation}
I write $\Omega = \omega - \omega_0$ for the distance from the carrier.

\textbf{Step 2.} I complete the square in the exponent:
\begin{align}-at^2 + i\Omega t &= -a\left(t^2 - \frac{i\Omega}{a}t\right)\\
&= -a\left[\left(t - \frac{i\Omega}{2a}\right)^2 - \left(\frac{i\Omega}{2a}\right)^2\right]\\
&= -a\left(t - \frac{i\Omega}{2a}\right)^2 + a\,\frac{i^2\Omega^2}{4a^2}\\
&= -a\left(t - \frac{i\Omega}{2a}\right)^2 - \frac{\Omega^2}{4a} .\end{align}

\textbf{Step 3.} The second term does not depend on $t$ and leaves the integral. The first is a Gaussian, and
$\int e^{-a(t-t_0)^2}dt = \sqrt{\pi/a}$ for any shift $t_0$, so
\begin{equation}\tilde E(\omega) = \sqrt{\frac{\pi}{a}}\,\exp\!\left[-\frac{(\omega-\omega_0)^2}{4a}\right] .\end{equation}
This is the spectrum of [Rom03], $E(\omega) \propto \exp[-(\omega-\omega_c)^2/4a]$.

\textbf{Step 4.} With $a = 2\ln 2/\tau^2$, $1/4a = \tau^2/(8\ln 2)$. Dropping the constant $\sqrt{\pi/a}$,
\begin{equation}a(\omega) = \exp\!\left[-\frac{(\omega-\omega_0)^2\tau^2}{8\ln 2}\right]
= \exp\!\left[-\frac{(\omega-\omega_0)^2}{2s^2}\right],\qquad s^2 = \frac{4\ln 2}{\tau^2},\quad
s = \frac{2\sqrt{\ln 2}}{\tau} .\label{eq:spec}\end{equation}
This is the function in \texttt{mcdo.pulsed2ce.spectrum}.

\textbf{Step 5, a check with the numbers of [Rom03].} The power spectrum is $|a(\omega)|^2 = e^{-\Omega^2/s^2}$. It is
one half at $\Omega^2 = s^2\ln 2$, so its FWHM is $\Delta\omega = 2s\sqrt{\ln 2} = 4\ln 2/\tau$. In frequency,
\begin{equation}\Delta\nu = \frac{\Delta\omega}{2\pi} = \frac{4\ln 2}{2\pi\tau} = \frac{0.441}{\tau} .\end{equation}
For $\tau = 1$ fs this is $4.4\times10^{14}$ Hz. With $\nu_c = 3.997\times10^{14}$ Hz the band runs from
$\nu_c - \Delta\nu/2 = 1.79\times10^{14}$ Hz (1.67 $\mu$m) to $\nu_c + \Delta\nu/2 = 6.20\times10^{14}$ Hz (483 nm),
the values quoted in [Rom03].

\section{The pulse at the focus, and the meaning of a delay}
\textbf{Step 1.} Each frequency of the pulse is a continuous wave, and the lens focuses it into the field
$\mathbf{E}_\omega(P)$ of the papers ([CS00] Eq.~1 for the scalar model, [Rom03] Eqs.~1--7 for the vector model).
The diffraction integral has the prefactor $1/\lambda \propto \omega$. The pulse at the point $P$ is the sum
\begin{equation}\mathbf{E}(P,t) = \int a(\omega)\,\frac{\omega}{\omega_0}\,\mathbf{E}_\omega(P)\,e^{-i\omega t}\,d\omega .\label{eq:field}\end{equation}

\textbf{Step 2, when the pulse arrives.} On the axis of a beam the focal field is a sum over the rays of the cone,
and the ray at the angle $\vartheta$ to the axis carries the phase $kz_b\cos\vartheta$, with $k = n\omega/c$ (Part~1 of
the derivations, Section~7). Near the focus the sum has the phase of the mean,
\begin{equation}\mathbf{E}_\omega(z_b) \propto e^{ikz_b\langle\cos\vartheta\rangle} ,\end{equation}
where $\langle\cos\vartheta\rangle$ is the mean of $\cos\vartheta$ over the cone, each ray weighted by its amplitude. In
the paraxial model it is $1 - (\mathrm{NA}/n)^2/4$, which is the phase $kz - u/4$ of Part~1, Section~3; in the vector
model the weight is $\sqrt t\,(1 + t)$, with $t = \cos\vartheta$ (Part~1, Section~7). Then
\begin{equation}e^{ikz_b\langle\cos\vartheta\rangle}\,e^{-i\omega t} = e^{-i\omega(t - nz_b\langle\cos\vartheta\rangle/c)} ,\end{equation}
and the sum over $\omega$ is the same function of $t - nz_b\langle\cos\vartheta\rangle/c$ that it is of $t$ at the
focus. So the pulse is at the focus at $t = 0$ and at the point $z_b$ at
\begin{equation}t_{\text{arrival}} = \frac{n z_b\langle\cos\vartheta\rangle}{c} = \frac{z_b}{v},\qquad
v = \frac{c}{n\langle\cos\vartheta\rangle} .\label{eq:arr}\end{equation}
The focal phase $-u/4$ grows in proportion to $\omega$, so it moves the pulse as well as the carrier: near the focus
the pulse travels along the axis at $v$, faster than $c/n$, by 0.25\% at NA 0.1, 7\% at NA 0.5 and 33\% at NA 0.9 in
air (vector model).

\textbf{A check with the engine.} The engine adds the frequencies with their focal fields, so it contains this speed.
For 20-fs pulses at 750 nm in air, the centroid in time of $|\mathbf{E}(z_b,t)|^2$ at $z_b = 0.25\ \mu$m is $nz_b/c$
times 0.9975, 0.9771, 0.9346, 0.8648 and 0.7541 at NA 0.1, 0.3, 0.5, 0.7 and 0.9 (vector model), where
$\langle\cos\vartheta\rangle$ is 0.9975, 0.9771, 0.9346, 0.8647 and 0.7517. In the paraxial model the centroid is
$nz_b(1 - \mathrm{NA}^2/4)/c$ at all $z_b$, since the rest of the field, $\mathrm{sinc}(u/4)$, is real. Farther from a
focus at NA 0.9 the centroid lies between $\langle\cos\vartheta\rangle$ and 1 (0.81 to 0.90 of $nz_b/c$ from 2 to
15~$\mu$m), where the 2CE signal is small.

\textbf{Step 3, the delay.} A delay line makes the whole field of beam 2 arrive later by $\Delta t$:
\begin{equation}\mathbf{E}_2(P,t) \longrightarrow \mathbf{E}_2(P,t-\Delta t)
= \int a_2(\omega)\,\frac{\omega}{\omega_2}\,\mathbf{E}_\omega(P)\,e^{-i\omega t}\,e^{+i\omega\Delta t}\,d\omega .\label{eq:delay}\end{equation}
Each frequency of beam 2 is multiplied by the phase $e^{i\omega\Delta t}$. This is what the code applies, with the
carrier included.

\section{The signal: photon pairs that reach the level}
\textbf{Step 1.} For a transition from the ground state $g$ to the level $f$ through intermediate states $m$,
second-order perturbation theory gives the amplitude
\begin{equation}a_f = \left(\frac{-i}{\hbar}\right)^2\sum_m\mu_{fm}\mu_{mg}
\int_{-\infty}^{\infty}dt\,E(t)\,e^{i\omega_{fm}t}\int_{-\infty}^{t}dt'\,E(t')\,e^{i\omega_{mg}t'} .\end{equation}

\textbf{Step 2.} The intermediate states are not resonant with the light: $\omega_{mg}$ differs from all frequencies
of the pulses by much more than their bandwidth and than the distance between the two carriers,
$\omega_2 - \omega_1 = 0.36$ rad/fs. For a component $E(t') \propto e^{-i\omega t'}$ the inner integral
is then
\begin{equation}\int_{-\infty}^{t}e^{i(\omega_{mg}-\omega)t'}dt' = \frac{e^{i(\omega_{mg}-\omega)t}}{i(\omega_{mg}-\omega)} ,\end{equation}
and since $\omega_{mg} - \omega$ is nearly the same for all $\omega$ of the two pulses, the inner integral is
proportional to $E(t)\,e^{i\omega_{mg}t}$.

\textbf{Step 3.} The two exponentials combine, $e^{i\omega_{fm}t}e^{i\omega_{mg}t} = e^{i\omega_{fg}t}$, and
$\omega_{fg} = \omega_e$. So, up to a constant,
\begin{equation}A(P) = \int_{-\infty}^{\infty}E^2(P,t)\,e^{i\omega_et}\,dt .\end{equation}
The level is taken to be narrow compared with the spectra, so one frequency $\omega_e$ enters.

\textbf{Step 4, the same amplitude in frequencies.} Writing each factor of $E^2$ through its spectrum,
$E(t) = \frac{1}{2\pi}\int\tilde E(\omega)e^{-i\omega t}d\omega$,
\begin{align}A &= \frac{1}{(2\pi)^2}\iint\tilde E(\omega_a)\tilde E(\omega_b)
\left[\int e^{i(\omega_e-\omega_a-\omega_b)t}dt\right]d\omega_a\,d\omega_b\\
&= \frac{1}{(2\pi)^2}\iint\tilde E(\omega_a)\tilde E(\omega_b)\;2\pi\,\delta(\omega_e-\omega_a-\omega_b)\;d\omega_a\,d\omega_b\\
&= \frac{1}{2\pi}\int\tilde E(\omega_a)\,\tilde E(\omega_e-\omega_a)\,d\omega_a .\label{eq:pairs}\end{align}
The delta function says that a pair of photons contributes when $\omega_a + \omega_b = \omega_e$, and the integral
says that all such pairs add with their phases.

\textbf{Step 5, the three kinds of pairs.} With two beams, $E = E_1 + E_2$ and
\begin{equation}E^2 = E_1^2 + 2E_1E_2 + E_2^2 ,\end{equation}
so
\begin{equation}A = A^{11} + 2A^{12} + A^{22},\qquad A^{ab} = \int E_a(P,t)\,E_b(P,t)\,e^{i\omega_et}\,dt .\end{equation}
\begin{center}\small\begin{tabular}{lll}\toprule
Term & Photons & Name\\\midrule
$2A^{12}$ & one from each beam, 750 + 656 nm & 350-nm pairs\\
$A^{11}$ & both from beam 1, 750 + 750 nm & 375-nm pairs\\
$A^{22}$ & both from beam 2, 656 + 656 nm & 328-nm pairs\\\bottomrule\end{tabular}\end{center}
The names give the wavelength that two photons at the carriers would reach. The 375- and 328-nm pairs reach the
350-nm level through the edges of the spectra. The signal is
\begin{equation}F(P) = |A(P)|^2 .\end{equation}

\textbf{Step 6, the limit of long pulses.} For long pulses each spectrum is a line, $E_a \to U_a e^{-i\omega_at}$,
with $U_a$ the continuous-wave field. Then
\begin{equation}A^{ab} \to U_aU_b\int e^{i(\omega_e-\omega_a-\omega_b)t}dt ,\end{equation}
which is zero unless $\omega_a + \omega_b = \omega_e$. For 2CE, $\omega_1 + \omega_2 = \omega_e$ but
$2\omega_1 \ne \omega_e$ and $2\omega_2 \ne \omega_e$, so the 350-nm pairs alone remain:
\begin{equation}F \to |2U_1U_2|^2 = 4|U_1|^2|U_2|^2 \propto I_1I_2 ,\end{equation}
which is Eq.~2 of [CS00]. For 2PE, $\omega_1 = \omega_2 = \omega_e/2$ and all three terms remain:
\begin{equation}F \to |U_1^2 + 2U_1U_2 + U_2^2|^2 = |U_1 + U_2|^4 ,\end{equation}
which is Eq.~3 of [CS00]. The pulsed model therefore contains the continuous-beam results of the papers as its
limit.

\section{How large the 375- and 328-nm pairs are}
\textbf{Step 1.} At one point, with both pulses arriving together, the amplitude of a pair of beams $a, b$ is the
integral of Step 4 with the two Gaussian spectra of Eq.~\ref{eq:spec}:
\begin{equation}A^{ab} \propto \int\exp\!\left[-\frac{(\omega-\omega_a)^2}{2s^2}\right]
\exp\!\left[-\frac{(\omega_e-\omega-\omega_b)^2}{2s^2}\right]d\omega .\end{equation}

\textbf{Step 2.} I write $x = \omega - \omega_a$ and $\Delta = \omega_e - \omega_a - \omega_b$, the distance of
the level from the sum of the two carriers. Then $\omega_e - \omega - \omega_b = \Delta - x$ and the exponent is
\begin{align}-\frac{1}{2s^2}\left[x^2 + (\Delta - x)^2\right]
&= -\frac{1}{2s^2}\left[2x^2 - 2\Delta x + \Delta^2\right]\\
&= -\frac{1}{2s^2}\left[2\left(x - \frac{\Delta}{2}\right)^2 - \frac{\Delta^2}{2} + \Delta^2\right]\\
&= -\frac{1}{s^2}\left(x - \frac{\Delta}{2}\right)^2 - \frac{\Delta^2}{4s^2} .\end{align}

\textbf{Step 3.} The integral over $x$ of the first term is $s\sqrt{\pi}$ for any $\Delta$. So
\begin{equation}A^{ab} \propto \exp\!\left[-\frac{\Delta^2}{4s^2}\right]
= \exp\!\left[-\frac{(\omega_e-\omega_a-\omega_b)^2\tau^2}{16\ln 2}\right],\end{equation}
using $4s^2 = 16\ln 2/\tau^2$.

\textbf{Step 4.} For the 350-nm pairs $\Delta = \omega_e - \omega_1 - \omega_2 \approx 0$ and the exponential is 1.
For the 375-nm pairs $\Delta = \omega_e - 2\omega_1 = 0.3589$ rad/fs. The first term also lacks the factor 2 of
$2A^{12}$. The factor $\omega/\omega_0$ of Eq.~\ref{eq:field} enters once for each photon; the 375-nm pairs that reach the level
have both photons near $\omega_e/2$, which gives $(\omega_e/2\omega_1)^2 = 1.148$. So
\begin{equation}\frac{|A^{11}|}{|2A^{12}|} = \frac12\left(\frac{\omega_e}{2\omega_1}\right)^2
\exp\!\left[-\frac{(\omega_e-2\omega_1)^2\tau^2}{16\ln 2}\right] ,\label{eq:ratio}\end{equation}
and the same with $\omega_2$ for the 328-nm pairs, where $\omega_e - 2\omega_2 = -0.3609$ rad/fs and
$(\omega_e/2\omega_2)^2 = 0.878$.

\textbf{Step 5, the numbers.} At $\tau = 10$ fs: $(0.3589)^2\times100/(16\ln 2) = 12.88/11.09 = 1.161$,
$e^{-1.161} = 0.313$, and $\tfrac12\times1.148\times0.313 = 0.180$.
\begin{center}\small\begin{tabular}{lcccc}\toprule
 & \multicolumn{2}{c}{375-nm pairs over 350-nm pairs} & \multicolumn{2}{c}{328-nm pairs over 350-nm pairs}\\
$\tau$ & Eq.~\ref{eq:ratio} & computed at the focus & Eq.~\ref{eq:ratio} & computed at the focus\\\midrule
5 fs & 0.429 & 0.430 & 0.327 & 0.327\\
10 fs & 0.180 & 0.180 & 0.136 & 0.136\\
20 fs & 0.0055 & 0.005 & 0.0040 & 0.004\\\bottomrule\end{tabular}\end{center}
[LS03] state that spectral broadening is negligible for pulse widths longer than 10 fs. Eq.~\ref{eq:ratio} gives the size:
below 1\% at 20 fs, 18\% at 10 fs.

\section{A delay at a fixed point: the temporal overlap}
\textbf{Step 1.} At a point where the two pulses arrive together when $\Delta t = 0$, the envelopes with a delay
are, from Eq.~\ref{eq:pulse},
\begin{equation}e_1(t) = \exp\!\left(-2\ln 2\,\frac{t^2}{\tau^2}\right),\qquad
e_2(t) = \exp\!\left(-2\ln 2\,\frac{(t-\Delta t)^2}{\tau^2}\right) .\end{equation}

\textbf{Step 2.} The 350-nm pairs are $2A^{12} \propto \int e_1(t)\,e_2(t)\,dt$ (the carriers cancel against
$e^{i\omega_et}$, since $\omega_1 + \omega_2 = \omega_e$). The product has the exponent
\begin{equation}-\frac{2\ln 2}{\tau^2}\left[t^2 + (t-\Delta t)^2\right] .\end{equation}

\textbf{Step 3.} I complete the square in the bracket:
\begin{align}t^2 + (t-\Delta t)^2 &= 2t^2 - 2t\,\Delta t + \Delta t^2\\
&= 2\left(t^2 - t\,\Delta t + \frac{\Delta t^2}{4}\right) - \frac{\Delta t^2}{2} + \Delta t^2\\
&= 2\left(t - \frac{\Delta t}{2}\right)^2 + \frac{\Delta t^2}{2} .\end{align}

\textbf{Step 4.} So
\begin{equation}e_1(t)\,e_2(t) = \exp\!\left[-\frac{4\ln 2}{\tau^2}\left(t - \frac{\Delta t}{2}\right)^2\right]
\exp\!\left(-\ln 2\,\frac{\Delta t^2}{\tau^2}\right) .\end{equation}
The first factor is a Gaussian centered halfway between the two pulses, and its integral over $t$ does not depend
on $\Delta t$. The second factor leaves the integral:
\begin{equation}|2A^{12}| \propto \exp\!\left(-\ln 2\,\frac{\Delta t^2}{\tau^2}\right) .\label{eq:g}\end{equation}

\textbf{Step 5.} The signal is the square,
\begin{equation}\frac{F(\Delta t)}{F(0)} = \exp\!\left(-2\ln 2\,\frac{\Delta t^2}{\tau^2}\right) .\label{eq:overlap}\end{equation}
\begin{center}\small\begin{tabular}{lcccc}\toprule
$\Delta t$ & $0.5\tau$ & $\tau/\sqrt2 = 0.71\tau$ & $\tau$ & $2\tau$\\\midrule
$F(\Delta t)/F(0)$ & $2^{-1/2} = 0.707$ & $2^{-1} = 0.5$ & $2^{-2} = 0.25$ & $2^{-8} = 0.0039$\\
stored, collinear lenses, 100 fs & 0.70712 & & 0.25001 & 0.00391\\
stored, collinear lenses, 5 fs & 0.71095 & & 0.25545 & 0.00425\\\bottomrule\end{tabular}\end{center}
Each value is a power of 2, since $\exp(-2\ln 2\,x^2) = 2^{-2x^2}$. The difference of 2\% at 5 fs comes from the
factor $\omega/\omega_0$ of Eq.~\ref{eq:field}, once for each photon. With $x = \omega - \omega_1$ for the photon of
beam 1, the photon of beam 2 has $\omega_e - \omega = \omega_2 - x$, and the two Gaussian spectra multiply to
$e^{-x^2/s^2}$. The amplitude of the 350-nm pairs at a delay is then
\begin{equation}|2A^{12}| \propto \left|\int e^{-x^2/s^2}\left(1 + \frac{x}{\omega_1}\right)\left(1 - \frac{x}{\omega_2}\right)
e^{-ix\Delta t}\,dx\right| ,\end{equation}
where the product of the two factors is
\begin{equation}\left(1 + \frac{x}{\omega_1}\right)\left(1 - \frac{x}{\omega_2}\right) = 1 + bx - \frac{x^2}{\omega_1\omega_2},\qquad
b = \frac{1}{\omega_1} - \frac{1}{\omega_2} .\end{equation}
With $G_0 = \int e^{-x^2/s^2}e^{-ix\Delta t}dx = s\sqrt\pi\,e^{-s^2\Delta t^2/4}$, the factors $x$ and $x^2$ become
$i\,\partial/\partial\Delta t$ applied once and twice, $-is^2\Delta t\,G_0/2$ and $(s^2/2 - s^4\Delta t^2/4)\,G_0$, so
\begin{equation}\int e^{-x^2/s^2}\left(1 + bx - \frac{x^2}{\omega_1\omega_2}\right)e^{-ix\Delta t}dx
= G_0\left[1 - \frac{s^2}{2\omega_1\omega_2} + \frac{s^4\Delta t^2}{4\omega_1\omega_2} - \frac{ibs^2\Delta t}{2}\right] .\end{equation}
The term in $x^2$ narrows the product of the spectra and the term in $x$ adds an imaginary part; both raise the
overlap at a delay, so the signal falls more slowly than Eq.~\ref{eq:overlap}. At 5 fs the square of the bracket,
over its value at $\Delta t = 0$, is 1.0054, 1.0218 and 1.0886 at $\Delta t = 0.5\tau$, $\tau$ and $2\tau$, which gives
0.71094, 0.25545 and 0.00425; the stored values are 0.71095, 0.25545 and 0.00425.

\textbf{Step 6.} The overlap of the two intensities gives the same function:
$\int I_1(t)I_2(t-\Delta t)\,dt \propto \exp(-2\ln 2\,\Delta t^2/\tau^2)$, by Steps 2 to 4 with the intensities
of Eq.~\ref{eq:pulse} in place of the envelopes. This is the sense in which [CS00] require $I_1$ and $I_2$ to be
``simultaneously present in both space and time''.

\section{Facing lenses: where the pulses meet}
\textbf{Step 1, arrival times.} With facing lenses, beam 1 travels along $+z$ and beam 2 along $-z$. From Eq.~\ref{eq:arr},
pulse 1 reaches the point $z$ at $t_1 = z/v$. For beam 2, $z_b = -z$, and it is delayed, so
$t_2 = -z/v + \Delta t$.

\textbf{Step 2.} At the point $z$ the two pulses are apart in time by
\begin{equation}t_1 - t_2 = \frac{2z}{v} - \Delta t .\end{equation}
Eq.~\ref{eq:overlap} holds at each point with this difference in place of $\Delta t$, and the continuous-beam intensities give
the strength of each pulse there:
\begin{equation}F_{2c}(z,\Delta t) = I_1(z)\,I_2(z)\,
\exp\!\left[-\frac{2\ln 2}{\tau^2}\left(\frac{2z}{v} - \Delta t\right)^2\right] .\label{eq:gate}\end{equation}

\textbf{Step 3, the meeting point.} The exponential is 1 where $2z/v = \Delta t$:
\begin{equation}z_m = \frac{v\,\Delta t}{2} = \frac{c\,\Delta t}{2n\langle\cos\vartheta\rangle} .\end{equation}
For $\Delta t = 100$ fs in air, $z_m = 0.29979\times100/(2\langle\cos\vartheta\rangle) = 15.0\ \mu\mathrm{m}/\langle\cos\vartheta\rangle$:
15.0~$\mu$m at NA 0.1 and 16.0~$\mu$m at NA 0.5.

\textbf{Step 4, the length of the stretch where the pulses meet.} With $2z/v - \Delta t = (2/v)(z - z_m)$, the
exponential is
\begin{equation}\exp\!\left[-\frac{8\ln 2}{v^2\tau^2}(z - z_m)^2\right]
= \exp\!\left[-4\ln 2\,\frac{(z-z_m)^2}{G^2}\right],\qquad G = \frac{v\tau}{\sqrt2} = \frac{c\tau}{\sqrt2\,n\langle\cos\vartheta\rangle} ,\end{equation}
since $4\ln 2/G^2 = 8\ln 2/v^2\tau^2$ gives $G^2 = v^2\tau^2/2$. By Eq.~\ref{eq:pulse} a Gaussian written this way has
FWHM $G$. In air at low NA, $G = 21.2\ \mu$m at 100 fs, 4.2 $\mu$m at 20 fs and 2.1 $\mu$m at 10 fs; at NA 0.5 each
is 7\% longer.

\textbf{Step 5, the peak when the focus is approximated by a Gaussian.} Near the focus I write the
continuous-beam product as a Gaussian of FWHM $D$, the continuous-beam depth of focus of 2CE:
$I_1I_2 \approx \exp(-4\ln 2\,z^2/D^2)$. Then
\begin{equation}F_{2c}(z,\Delta t) \approx \exp\!\left[-4\ln 2\left(\frac{z^2}{D^2} + \frac{(z-z_m)^2}{G^2}\right)\right] .\end{equation}
The peak is where the derivative of the bracket is zero:
\begin{align}\frac{2z}{D^2} + \frac{2(z - z_m)}{G^2} &= 0\\
z\left(\frac{1}{D^2} + \frac{1}{G^2}\right) &= \frac{z_m}{G^2}\\
z_p &= z_m\,\frac{D^2}{D^2 + G^2} .\label{eq:zp}\end{align}

\textbf{Step 6, the height of the peak.} At $z_p$, $z_p - z_m = -z_mG^2/(D^2+G^2)$, so
\begin{align}\frac{z_p^2}{D^2} + \frac{(z_p - z_m)^2}{G^2}
&= \frac{z_m^2D^2}{(D^2+G^2)^2} + \frac{z_m^2G^2}{(D^2+G^2)^2}
= \frac{z_m^2}{D^2+G^2} ,\end{align}
and
\begin{equation}\frac{F_{\text{peak}}(\Delta t)}{F_{\text{peak}}(0)}
= \exp\!\left[-4\ln 2\,\frac{z_m^2}{D^2 + G^2}\right] .\end{equation}

\textbf{Step 7, the width.} The sum of the two exponents has the coefficient $1/D^2 + 1/G^2 = 1/W^2$, so the
profile along the axis has FWHM
\begin{equation}W = \frac{DG}{\sqrt{D^2 + G^2}} ,\label{eq:W}\end{equation}
which is $G$ where the focus is long ($D \gg G$) and $D$ where it is short.

\textbf{Step 8, the two limits at $\Delta t = \tau$.} Here $z_m = v\tau/2 = G/\sqrt2$, so $z_m^2 = G^2/2$ and
\begin{equation}\frac{F_{\text{peak}}(\tau)}{F_{\text{peak}}(0)} = \exp\!\left[-2\ln 2\,\frac{G^2}{D^2 + G^2}\right] .\label{eq:peaktau}\end{equation}
For a short focus, $D \ll G$, this is $e^{-2\ln 2} = 0.25$ and $z_p \to 0$: the peak stays at the focus and falls
with the temporal overlap. For a long focus, $D \gg G$, it is 1 and $z_p \to z_m$: the peak moves with the
meeting point and is kept.

\textbf{Step 9, comparison with the stored results.} For 100 fs in air, $G = 21.2\ \mu\mathrm{m}/\langle\cos\vartheta\rangle$
and $z_m = 15.0\ \mu\mathrm{m}/\langle\cos\vartheta\rangle$. The stored values are from the vector model, so $D$ and
$\langle\cos\vartheta\rangle$ are those of the vector model too.
\begin{center}\small\begin{tabular}{lcccccccc}\toprule
 & & & \multicolumn{2}{c}{peak at $\Delta t = \tau$} & \multicolumn{2}{c}{position ($\mu$m)} &
\multicolumn{2}{c}{FWHM, no delay ($\mu$m)}\\
NA & $D$ ($\mu$m) & $\langle\cos\vartheta\rangle$ & Eq.~\ref{eq:peaktau} & stored & Eq.~\ref{eq:zp} & stored & Eq.~\ref{eq:W} & stored\\\midrule
0.1 & 88.8 & 0.9975 & 0.928 & 0.930 & 14.21 & 14.24 & 20.67 & 20.69\\
0.2 & 22.0 & 0.990 & 0.510 & 0.513 & 7.79 & 7.78 & 15.36 & 15.43\\
0.3 & 9.67 & 0.977 & 0.315 & 0.316 & 2.54 & 2.58 & 8.83 & 8.85\\
0.5 & 3.33 & 0.935 & 0.257 & 0.258 & 0.34 & 0.34 & 3.29 & 3.29\\\bottomrule\end{tabular}\end{center}
With pulses at $c/n$ and the depth of focus of the scalar model, as on the slides, the formula gives a peak of 0.928,
0.517, 0.320 and 0.260 and, at NA 0.5, a position of 0.41~$\mu$m and an FWHM of 3.51~$\mu$m. Without the Gaussian
approximation, Eq.~\ref{eq:gate} with the computed $I_1I_2$ agrees with the stored profiles to 0.002 of the peak at
NA 0.1 to 0.9 and at 10, 20 and 100 fs, in both models. With pulses at $c/n$ the differences are 0.004 at 100~fs,
0.019 at 20~fs and 0.039 at 10~fs, so the stored profiles contain the speed $v$.
\begin{center}\includegraphics[width=0.95\textwidth]{../../output/08_two_color/cambaliza2000/pulses/delay_sweep/delay_facing_vs_na.png}\end{center}

\section{Fringes against the delay}
\textbf{Step 1.} From Eq.~\ref{eq:delay}, the delay multiplies each photon of beam 2 by $e^{i\omega\Delta t}$. I count the
photons of beam 2 in each kind of pair.
\begin{itemize}
\item A 375-nm pair has no photon of beam 2. $A^{11}$ does not change.
\item A 350-nm pair has one, of frequency $\omega_b$ near $\omega_2$. $2A^{12}$ gains about $e^{i\omega_2\Delta t}$,
and its size falls as $g(\Delta t) = \exp(-\ln 2\,\Delta t^2/\tau^2)$ (Eq.~\ref{eq:g}).
\item A 328-nm pair has two, with $\omega_a + \omega_b = \omega_e$ exactly (Eq.~\ref{eq:pairs}). $A^{22}$ gains
$e^{i\omega_a\Delta t}e^{i\omega_b\Delta t} = e^{i\omega_e\Delta t}$, and its size does not change.
\end{itemize}
\begin{equation}A(\Delta t) = A^{11} + 2A^{12}_0\,g(\Delta t)\,e^{i\omega_2\Delta t} + A^{22}\,e^{i\omega_e\Delta t} .\end{equation}

\textbf{Step 2.} $|A|^2$ contains the three cross terms of this sum. Each oscillates with the difference of the
phases of its two terms, with period $2\pi$ over that difference:
\begin{center}\small\begin{tabular}{llll}\toprule
Pairs that interfere & Phase difference & Period & From the computed phases\\\midrule
350-nm and 375-nm & $\omega_2\Delta t$ & $2\pi/\omega_2 = \lambda_2/c = 2.19$ fs & 2.19 fs\\
350-nm and 328-nm & $(\omega_e - \omega_2)\Delta t = \omega_1\Delta t$ & $\lambda_1/c = 2.50$ fs & 2.50 fs\\
375-nm and 328-nm & $\omega_e\Delta t$ & $\lambda_e/c = 1.17$ fs & 1.17 fs\\\bottomrule\end{tabular}\end{center}
The first two fall with $g(\Delta t)$. The third keeps its depth at all delays, since neither of its terms needs
the two pulses to overlap.

\textbf{Step 3, when there are no fringes.} The fringes need the phase between the two beams to be fixed from
pulse to pulse. With two lasers that are not phase-locked the cross terms average to zero, and the signal of the
350-nm pairs, Eq.~\ref{eq:overlap}, remains. The fluorescence probes are also chosen against two-photon excitation by
$\lambda_1$ or $\lambda_2$ alone [LS03]. For both reasons the 2CE signal in the delay results is the 350-nm pairs,
and the total with all pairs is shown beside it below 20 fs.

\textbf{Step 4, the fringes in space.} Each term also carries the carriers of its beams along the axis. With
facing lenses beam 1 has $e^{ik_1z}$ and beam 2 has $e^{-ik_2z}$:
\begin{equation}A^{11} \propto e^{2ik_1z},\qquad 2A^{12} \propto e^{i(k_1-k_2)z},\qquad A^{22} \propto e^{-2ik_2z} .\end{equation}
The phase difference between $A^{11}$ and $2A^{12}$ is $(2k_1 - k_1 + k_2)z = (k_1+k_2)z$, and between $2A^{12}$
and $A^{22}$ it is also $(k_1+k_2)z$. The period along $z$ is
\begin{equation}\frac{2\pi}{k_1 + k_2} = \frac{1}{n(1/\lambda_1 + 1/\lambda_2)} = \frac{\lambda_e}{n} ,\end{equation}
350 nm in air. Near the focus each $k$ is multiplied by $\langle\cos\vartheta\rangle$ (Section~3, Step~2), so the fringes
are $\lambda_e/(n\langle\cos\vartheta\rangle)$ apart, slightly more than 350 nm at high NA, as for the 2PE standing wave
(Part~3 of the derivations, Section~3, Step~5).

\textbf{Step 5, why the depth depends on NA.} A delay changes the phases of Step 1, which moves the fringes of
Step 4 along $z$. The peak along the axis is the largest value on the line, so it follows a fringe as it moves.
Where the focus is many fringes long, some fringe is always near the top of the envelope and the peak changes
little. Where the focus is about one fringe long (high NA), the peak rises and falls with the fringe. With
collinear lenses the three terms have the same carrier, $e^{i(k_a+k_b)z}$ with $k_a + k_b = n\omega_e/c$ for all of
them, so there are no fringes in space and the depth against the delay is the same at all NAs.
\begin{center}\small\begin{tabular}{lcccc}\toprule
Largest gap, all pairs minus 350-nm pairs & NA 0.1 & 0.3 & 0.5 & 0.7\\\midrule
collinear, 5 fs / 10 fs / 20 fs & \multicolumn{4}{c}{0.917 / 0.710 / 0.036 at all NAs}\\
crossed, along $x$, 5 fs & 0.062 & 0.197 & 0.416 & 0.652\\
facing, along $z$, 5 fs & 0.014 & 0.013 & 0.033 & 0.070\\\bottomrule\end{tabular}\end{center}
With crossed lenses the fringes run along the diagonal of the $x$--$z$ plane. Over the whole plane the gap at
5 fs is 0.278 at NA 0.5 and 0.478 at NA 0.7, about two thirds of the values along $x$.

\section{The longitudinal component of the vector field}
The factor of $E_z$, which the two-beam code had with an extra $i$ until 2026-11-10, is derived step by step in
Part 2 of the derivations, \texttt{derivations/2\_vector\_focal\_field.pdf}: Section 3 derives the three components
and the constant, Section 4 checks the factor with Gauss's law, and Section 7 states what the factor changes. In short,
$E_z = -2I_1\cos\psi$ is real; $|E_z|^2$, the intensity of one beam and the 2CE signal are the same with either factor;
of the signals, the factor changes 2PE with the polarization in the plane of the two beams, which the pulse and delay runs do not use.

\section*{References}
\small
[CS00] M. O. Cambaliza and C. Saloma, Opt.\ Commun.\ 184, 25--35 (2000).\\
\relax[LS02] M. Lim and C. Saloma, Opt.\ Commun.\ 207, 121--130 (2002).\\
\relax[LS03] M. Lim and C. Saloma, Appl.\ Opt.\ 42(17), 3398--3406 (2003).\\
\relax[Rom03] K. M. Romallosa, J. Bantang and C. Saloma, Phys.\ Rev.\ A 68, 033812 (2003).\\
B. Richards and E. Wolf, Proc.\ R.\ Soc.\ London Ser.\ A 253, 358--379 (1959).
\end{document}
